% ===================================================================
\section{Introduction}
\label{sec:intro}

Since 28~February 2026, shipping through the Strait of Hormuz has been largely halted amid
the 2026 Iran war, the chokepoint sitting at the centre of the conflict
\citep{worldbank2026}. The strait is the narrowest binding constraint in the global
oil system, with roughly one-fifth of world petroleum-liquids consumption
\citep{eia_hormuz2025} and about 27\% of global seaborne oil trade \citep{crs_hormuz2026}
transiting the passage. What makes the
chokepoint so important is the absence of a redundant route. Hormuz is a narrow sea lane
between Iran and Oman with no maritime bypass, and the overland pipelines that skirt it
carry only a small share of the flow \citep{eia_hormuz2025}, so most Gulf crude has nowhere
else to go. Even
through recent regional conflict, traffic had never been halted, so the 2026 disruption is
the first time the strait has actually closed.

Crude oil is an input to almost every industry and a cost borne across nearly all
economies, so a disruption on this scale was always going to move prices. The shock
registered in Brent, the benchmark against which most internationally traded crude is
settled, which rose sharply once the lane was blocked and passed \$100 a barrel for the
first time in four years. This is not without precedent. The 2022 Russian invasion of
Ukraine drove Brent from about \$95 to \$130 within two weeks, the closest comparable
oil-supply shock of the past decade and the benchmark against which we later validate the
design.

The size and persistence of the move are what policy turns on. Emergency
strategic-reserve releases are sized against an implied price displacement, and the
International Energy Agency (IEA) coordinated the drawdown of some 180~million barrels
after the 2022 invasion \citep{iea2022release}. The monetary-policy choice to look through
an oil shock or to tighten depends on how large and how lasting it is. Fiscal
stabilisation in subsidising oil importers scales with the move. Sovereign hedging
programmes price disruption scenarios against exactly such magnitudes, the largest being
Mexico's annual Hacienda hedge \citep{jfi2023mexico}. The 2026 disruption has accordingly
drawn rapid analysis from the Kiel Institute, the Oxford Institute for Energy Studies, the
World Bank, and the IEA \citep{kiel2026,oies2026,worldbank2026,iea2026}.

Yet that policy-relevant effect is not the raw price jump. The jump conflates the disruption
with everything else moving markets at the same time, so its size is not the disruption's
alone, and as a single figure it says nothing about how the effect evolves. We therefore
pose a causal question.

\begin{quote}
What is the causal effect of the 2026 Strait of Hormuz disruption on the price of Brent
crude over a policy-relevant post-event window?
\end{quote}

We target the total causal effect rather than any single transmission channel, and we
estimate it as a path over the post-event window rather than a single number. Each policy
lever above acts on the trajectory of the effect, not on a peak-day figure, so the object
we need is the average gap between observed Brent and the price that would have prevailed
absent the disruption, traced over the whole window. Recovering that counterfactual path,
which has no directly observable analogue, is the methodological core of this paper.

Standard observational designs are ill-suited to it. Brent is a single, globally traded
benchmark, so difference-in-differences has no untreated control sharing its trend, and
instrumental variables has no credible instrument for a one-off geopolitical event. We use the synthetic control method (SCM) \citep{abadie2010,abadie2021}, which
forms a weighted basket of donor series whose pre-event co-movement with the outcome is
high and projects that basket forward as the counterfactual. SCM yields a full
counterfactual path, makes its identifying assumptions auditable through the donor pool,
and takes no structural stance on whether the move is supply, risk premium, or
expectations, since its estimand is a reduced-form total effect. The institutional analyses
cited before reach a price path by the opposite route, imposing an assumed barrel shortfall and
propagating it through a structural model. Ours reads the effect from the market and assumes no shortfall,
providing a complementary perspective.

This combination, SCM on a crude benchmark with cross-asset donors, has no direct
precedent. The closest price application uses cross-country donors for a regulated retail
market \citep{becker2021}, neither of which is available here. Because a single estimate
could be an artefact of the donor pool or the modelling choices, we set the Hormuz estimate
beside the 2022 Russian invasion, whose price move is already documented. Russia runs on the
same machinery but is estimated independently on its own pre-window, and the design
recovering its known move is what lends the Hormuz number its credibility. Thus, our contribution is twofold, a
data-driven average treatment effect on the treated (ATT) for the Hormuz disruption, and a
transparent template for validating cross-asset synthetic control under a single-unit
geopolitical treatment.

The remainder of the paper reviews the relevant literature (Section~\ref{sec:lit}),
sets out the factor model, estimand, and identifying assumptions
(Section~\ref{sec:theory}), describes the data, the estimators, and the validation
battery (Section~\ref{sec:data}), audits the donor pool donor by donor
(Section~\ref{sec:donoraudit}), and reports results for both events
(Section~\ref{sec:results}) before concluding (Section~\ref{sec:conclusion}).
